\subsection{A first-extension witness for the pointed Fermat quartic in characteristic five}

\begin{claim}{B.QM1}{Proved}{A constant cyclic thirteen-cover over \(\mathbf F_{125}\)}
Let
\[
 C/\mathbf F_5:\quad W^4=X^3Z-XZ^3,\qquad P=(1:0:0).
\]
Over $K=\mathbf F_{125}$ there is a constant cyclic, $P$-split,
\emph{\'etale} cover $D\to C$ of degree thirteen with
\[
  g_D=27,\qquad
  \frac{\#J_D(K)}{\#J_C(K)}\equiv1\pmod5.
\]
Its new Jacobian part is ordinary of dimension twenty-four.
Under an extension of $K$ of degree $n$, its relative class
number stays prime to five if and only if $4\nmid n$.
\end{claim}
\begin{proof}
All the following defining polynomials lie in $\mathbf F_5[x]$:
\[
\begin{split}
 h={}&x^6+x^2+4,\\
 N={}&x^{13}+3x^{11}+x^9+3x^7+2x^5+x,\\
 g={}&x^{18}+2x^{16}+x^{14}+3x^{12}+3x^{10}+2x^2+4.
\end{split}                                                    \tag{B.QM.1}
\]
Direct polynomial multiplication gives
\[
 g=N'h-2Nh',\qquad
       (x^3-x)g^2=N^3-Nh^4.                                    \tag{B.QM.2}
\]
The polynomials $x^3-x,g,h$ are squarefree and pairwise coprime,
and $N$ is coprime to $h$. Consequently
\[
 E:\ y^2=x^3-x\longrightarrow E,\qquad
       (x,y)\longmapsto
       \left(\frac{N(x)}{4h(x)^2},\frac{y g(x)}{2h(x)^3}\right)
                                                                  \tag{B.QM.3}
\]
is a separable degree-thirteen isogeny. Indeed (B.QM.2) verifies
the target equation, the map takes infinity to infinity, and
the rational function defining its first coordinate has degree
thirteen. A nonconstant morphism between elliptic curves taking
the origins to each other is an isogeny; its degree is thirteen
because the $x$-coordinate maps on source and target both have
degree two.

The kernel is constant over the displayed cubic field. A fully
explicit presentation and generator are
\[
 K=\mathbf F_5[z]/(z^3+4z^2+3z+1),\qquad
                    T=(z,z+3)\in E(K).                           \tag{B.QM.4}
\]
The cubic has no root in $\mathbf F_5$, so this is a field.
The elliptic addition law gives $13T=0$, with $T\ne0$.
The twelve nonzero multiples of $T$ have precisely the six
distinct roots of $h$ as their $x$-coordinates. These assertions
are also a short exact arithmetic certificate: its thirteen
points, including the identity, are listed in
\path{code/criterion_quartic13_verify.json}, in the basis
$1,z,z^2$. Thus they constitute the entire kernel of (B.QM.3).

Pull back (B.QM.3) along $C\to E$, $(X,W)\mapsto(X,W^2)$.
The normalization of the resulting connected curve is
\[
 D:\quad w^4=F(x),\qquad
 F(x)=4(x^3-x)g(x)^2h(x)^2,\qquad
 (X_C,W_C)=\left(\frac{N(x)}{4h(x)^2},\frac{w}{h(x)^2}\right).
                                                                  \tag{B.QM.5}
\]
One can verify connectedness directly: $F$ has simple zeros
at the three roots of $x^3-x$, so the fourth-root extension
has degree four even over an algebraic closure. Alternatively,
the degree-two extension $C/E$ and the odd-degree isogeny are
linearly disjoint. The map to $C$ is finite \emph{\'etale} by
base change of the isogeny; it is a constant cyclic cover over
$K$ by (B.QM.4). The fibre of $P$ is the rational kernel of
(B.QM.3), so it splits completely. Riemann--Hurwitz gives
$g_D=1+13(3-1)=27$.

We give the cohomological calculation explicitly enough to
reproduce it without an isogeny-search program. Here $\deg F=51$
and $\deg(gh)=24$. A basis of regular differentials is
\[
 \left\{\frac{x^jdx}{w}:0\le j\le11\right\},\qquad
 \left\{\frac{gh\,dx}{w^2}\right\},\qquad
 \left\{\frac{ghx^jdx}{w^3}:0\le j\le13\right\}.                \tag{B.QM.6}
\]
At a simple zero of $F$ the ramification index is four, while
at a double zero it is two. The factors $gh$ exactly remove
the poles in characters two and three at the latter zeros.
At infinity the orders of $x,w,dx$ are $-4,-51,-5$;
these give exactly the ranges in (B.QM.6). The forms are
independent by their characters and their powers of $x$, and
their number is the genus, proving that this is a full basis.

Put $h_1=1$ and $h_2=h_3=gh$. Since $5\equiv1\pmod4$,
Cartier in character $s$ is computed by
\[
 \mathcal C\!\left(\frac{h_s x^jdx}{w^s}\right)
   =\frac1{w^s}\,
      \mathcal C_{\mathbf P^1}(h_sx^jF^s dx),\qquad
 \mathcal C_{\mathbf P^1}\!\left(\sum a_i x^i dx\right)
        =\sum a_{5b+4}x^b dx.                                  \tag{B.QM.7}
\]
No coefficient roots need to be represented here because
all coefficients lie in $\mathbf F_5$. Divide the extracted
polynomial by $h_s$ and expand in the relevant part of
(B.QM.6); every division has zero remainder. This gives a
$27\times27$ matrix $M$ over $\mathbf F_5$. Its independently
computed characteristic polynomials are
\[
\begin{split}
 \det(U-M)
   &=(U-2)^5(U+2)^6
       (U^2-2U-1)^4(U^2+2U-1)^4,\\
 \det(U-M^3)&=(U-2)^{14}(U+2)^{13}.
\end{split}                                                    \tag{B.QM.8}
\]
For the old curve, the basis $dX/W^2,dX/W^3,X\,dX/W^3$ gives
\[
 M_C=\begin{pmatrix}3&0&0\\0&0&4\\0&1&0\end{pmatrix},\qquad
                \det(U-M_C^3)=(U-2)^2(U+2).                    \tag{B.QM.9}
\]
The entire matrix, its cube, and all intermediate polynomials
are in the certificate. The independent checker
\path{code/criterion_quartic13_verify.py} implements its own
prime-field polynomial arithmetic, the cubic field, the
elliptic group law, and (B.QM.7); it then computes the
characteristic polynomials by integer linear algebra reduced
modulo five. It imports no search-producer code.

Pullback and norm split old and new five-primary cohomology,
since thirteen is invertible there. Cartier and coherent
Frobenius are dual after a full finite-field cycle, and the
Artin--Schreier sequence identifies the stable part with
geometric five-torsion cohomology. Polynomial division of
(B.QM.8) by (B.QM.9) gives the new factor
\[
                  (U-2)^{12}(U+2)^{12}.                        \tag{B.QM.10}
\]
It has degree twenty-four and no zero root, proving the
claimed ordinarity. Its value at one is
$(-1)^{12}3^{12}=1$ in $\mathbf F_5$, so the relative
class number is a unit, with the asserted residue.
Both $2$ and $-2$ have multiplicative order four in
$\mathbf F_5$. The eigenvalues of a further $n$th Frobenius
power contain one exactly when $4\mid n$; possible Jordan
blocks do not affect this criterion. This proves the final
assertion.
\end{proof}

\begin{claim}{B.QM2}{Proved}{The exact first-field minimum and the pointed Fermat identification}
For the pointed quartic of \statementref{B.QM1}, the least
constant extension degree over which an odd-prime constant
cyclic $P$-split class-number witness exists is exactly three.
At that first field, thirteen is the least possible prime,
and \statementref{B.QM1} attains it. The same assertions hold
for the pointed Fermat quartic
\[
       U^4+V^4=R^4\quad\text{over }\mathbf F_5,
                   \qquad (U:V:R)=(1:0:1).
\]
This is a field-first minimum. The least prime when arbitrary
constant extension degrees are allowed is not determined here.
\end{claim}
\begin{proof}
Direct enumeration gives
\[
 \#C(\mathbf F_5)=8,\qquad
 \#C(\mathbf F_{25})=44,\qquad
 \#C(\mathbf F_{125})=104.
\]
Newton identities and the reciprocal functional equation
give the full Jacobian Frobenius polynomial
\[
          H_C(U)=(U^2+2U+5)^2(U^2-2U+5).                       \tag{B.QM.11}
\]
The independent certificate enumerates these points, verifies
the first three Newton identities, and checks the resulting
polynomial. Its consequences for the first three fields are
\[
\begin{split}
 \#J_C(\mathbf F_5)&=8^2\cdot4=256,\\
 \#J_C(\mathbf F_{25})&=32^3=32768,\\
 \#J_C(\mathbf F_{125})&=104^2\cdot148
                         =2^8\cdot13^2\cdot37.
\end{split}                                                    \tag{B.QM.12}
\]
A constant cyclic $P$-split cover of odd prime degree
$\ell\ne5$ gives a quotient $J_C(\mathbf F_{5^m})\twoheadrightarrow
\mathbf Z/\ell$, by the pointed Abel and Lang construction.
The first two orders in (B.QM.12) exclude every such odd prime
over degrees one and two. The third order permits only
thirteen and thirty-seven. The thirteen-cover already proved
therefore establishes both claimed minima.

Finally the characteristic-five identity
\[
       (X+2Z)^4-(X-2Z)^4=X^3Z-XZ^3
\]
shows that the invertible projective linear map
\[
       (X:W:Z)\longmapsto(X-2Z:W:X+2Z)
\]
carries $C$ to the displayed Fermat curve and $P$ to
$(1:0:1)$. Thus the result concerns that exact pointed
Fermat example, not merely another curve with the same
Jacobian point counts.
\end{proof}
